\chapter{PLC Data Type Reference (Current)}
\label{app:data-types}
This appendix lists the PLC data types (DDTs) used in EFC's Control Expert projects. Section~\ref{sec:udt} explains how they are used. The first three tables are the minimum member lists defined by this standard; the remaining tables document the types in the current program library, taken from the Control Expert data editor. Descriptions marked \tbd\ have not been confirmed yet.
\future{A replacement library (PCLib) has been proposed and is described in Appendix~\ref{app:data-types-future}. Until it is adopted, the types in this appendix remain the standard for existing and new projects.}

\note{Naming in the current library is not fully consistent: \texttt{AI\_Sensors} has no \texttt{DDT\_} prefix, and member capitalization differs between types (for example \texttt{Auto\_Perm} in \texttt{DDT\_Valve} and \texttt{AUTO\_PERM} in \texttt{DDT\_Motor\_VFD}). Align these the next time the library is revised.}

\section{DDT\_Process\_Value}
\begin{longtable}{@{}P{3.90cm}P{3.90cm}P{9.10cm}@{}}
\caption{DDT\_Process\_Value Minimum Members}\label{tab:ddt-process-value}\\
\toprule \rowcolor{tblhead}\textbf{Member} & \textbf{Type} & \textbf{Description}\\ \midrule
\endfirsthead
\toprule \rowcolor{tblhead}\textbf{Member} & \textbf{Type} & \textbf{Description}\\ \midrule
\endhead
Raw & INT or DINT & Unscaled channel value.\\
Raw\_Min / Raw\_Max & DINT & Configured raw endpoints.\\
PV & REAL & Scaled engineering value.\\
Eng\_Min / Eng\_Max & REAL & Engineering range endpoints.\\
Units & STRING[16] & Engineering-unit text or enumerated unit identifier.\\
Quality & DINT & Valid, uncertain, bad, stale, simulated, or platform equivalent.\\
Timestamp & DDT or native time & Time of last valid update where required.\\
LL / L / H / HH & DDT\_AI\_Alarm & Analog alarm instances.\\
HW\_Fault & BOOL & Module or channel diagnostic fault.\\
Underrange / Overrange & BOOL & Signal is below/above valid instrument range.\\
Simulated & BOOL & Value is under authorized simulation.\\
\bottomrule
\end{longtable}


\section{DDT\_AI\_Alarm}
\begin{longtable}{@{}P{3.90cm}P{3.90cm}P{9.10cm}@{}}
\caption{DDT\_AI\_Alarm Minimum Members}\label{tab:ddt-ai-alarm}\\
\toprule \rowcolor{tblhead}\textbf{Member} & \textbf{Type} & \textbf{Description}\\ \midrule
\endfirsthead
\toprule \rowcolor{tblhead}\textbf{Member} & \textbf{Type} & \textbf{Description}\\ \midrule
\endhead
SP & REAL & Rationalized alarm setpoint.\\
DB & REAL & Deadband in engineering units.\\
OnDelay / OffDelay & TIME & Qualified activation and clearing delays.\\
Enabled & BOOL & Alarm is configured for use.\\
Active & BOOL & Qualified alarm condition exists.\\
Latched & BOOL & Condition has latched where required.\\
Acked & BOOL & Acknowledgment state when managed in PLC.\\
Suppressed & BOOL & State-based designed suppression is active.\\
Shelved & BOOL & HMI/SCADA shelving indication when exchanged.\\
Priority & INT & Rationalized priority enumeration.\\
\bottomrule
\end{longtable}

\section{DDT\_Device}

\begin{longtable}{@{}P{4.29cm}P{3.90cm}P{8.71cm}@{}}
\caption{DDT\_Device Minimum Members}\label{tab:ddt-device}\\
\toprule \rowcolor{tblhead}\textbf{Member} & \textbf{Type} & \textbf{Description}\\ \midrule
\endfirsthead
\toprule \rowcolor{tblhead}\textbf{Member} & \textbf{Type} & \textbf{Description}\\ \midrule
\endhead
ModeReq / ModeAct & INT or enum & Requested and active operating mode.\\
StartReq / StopReq / ResetReq & BOOL & Momentary validated requests.\\
RunCmd & BOOL & Final logical run command before output mapping.\\
Running / Stopped & BOOL & Processed equipment status.\\
AutoPerm / ManPerm & BOOL & Mode-specific permissive summaries.\\
Interlock & DDT\_Interlock & Protective/control interlock summary.\\
FeedbackFault & BOOL & Command and feedback disagree after timeout.\\
Fault / Warning & BOOL & Aggregated condition summaries.\\
Available & BOOL & Equipment is healthy and eligible for control.\\
RuntimeSec & UDINT & Retentive accumulated runtime with controlled reset.\\
CommOK / DataValid & BOOL & Communications and exchanged data are valid.\\
Simulated & BOOL & Authorized simulation is active.\\
\bottomrule
\end{longtable}


\section{AI\_Sensors}
Analog sensor data with its four alarm states.

\begin{longtable}{@{}P{4.54cm}P{2.59cm}P{9.79cm}@{}}
\caption{AI\_Sensors Members}\label{tab:ddt-ai-sensors}\\
\toprule \rowcolor{tblhead}\textbf{Member} & \textbf{Type} & \textbf{Description}\\ \midrule
\endfirsthead
\multicolumn{3}{@{}l}{\small\textit{Table \thetable{} continued}} \\
\toprule \rowcolor{tblhead}\textbf{Member} & \textbf{Type} & \textbf{Description}\\ \midrule
\endhead
\texttt{PV} & REAL & Scaled process value in engineering units.\\
\texttt{Raw} & REAL & Unscaled value from the input channel.\\
\texttt{Rate} & REAL & Rate of change of the process value \tbd (units per time).\\
\texttt{Fraction} & REAL & \tbd\\
\texttt{LL\_Alarm} & BOOL & The low-low alarm is active.\\
\texttt{LL\_Alarm\_Acked} & BOOL & The low-low alarm has been acknowledged.\\
\texttt{Lo\_Alarm} & BOOL & The low alarm is active.\\
\texttt{Lo\_Alarm\_Acked} & BOOL & The low alarm has been acknowledged.\\
\texttt{Hi\_Alarm} & BOOL & The high alarm is active.\\
\texttt{Hi\_Alarm\_Acked} & BOOL & The high alarm has been acknowledged.\\
\texttt{HH\_Alarm} & BOOL & The high-high alarm is active.\\
\texttt{HH\_Alarm\_Acked} & BOOL & The high-high alarm has been acknowledged.\\
\bottomrule
\end{longtable}
\section{DDT\_Discrete}
Discrete input with an override for testing and maintenance.

\begin{longtable}{@{}P{4.54cm}P{2.59cm}P{9.79cm}@{}}
\caption{DDT\_Discrete Members}\label{tab:ddt-discrete}\\
\toprule \rowcolor{tblhead}\textbf{Member} & \textbf{Type} & \textbf{Description}\\ \midrule
\endfirsthead
\multicolumn{3}{@{}l}{\small\textit{Table \thetable{} continued}} \\
\toprule \rowcolor{tblhead}\textbf{Member} & \textbf{Type} & \textbf{Description}\\ \midrule
\endhead
\texttt{Override\_Enb} & BOOL & Override is enabled for this input; the PLC uses \texttt{Override\_Sig} instead of the field signal.\\
\texttt{Override\_Sig} & BOOL & Value used for the input while the override is enabled.\\
\texttt{Fault} & BOOL & Input fault (for example a channel or wiring fault) \tbd.\\
\bottomrule
\end{longtable}
\section{DDT\_Setpoint}
Operator setpoint with a default value.

\begin{longtable}{@{}P{4.54cm}P{2.59cm}P{9.79cm}@{}}
\caption{DDT\_Setpoint Members}\label{tab:ddt-setpoint}\\
\toprule \rowcolor{tblhead}\textbf{Member} & \textbf{Type} & \textbf{Description}\\ \midrule
\endfirsthead
\multicolumn{3}{@{}l}{\small\textit{Table \thetable{} continued}} \\
\toprule \rowcolor{tblhead}\textbf{Member} & \textbf{Type} & \textbf{Description}\\ \midrule
\endhead
\texttt{SP} & REAL & Active setpoint value.\\
\texttt{Default} & REAL & Default value loaded on first scan or reset \tbd.\\
\bottomrule
\end{longtable}
\section{DDT\_Valve}
On/off valve with limit switch feedback; used with \texttt{DFB\_Valve}.

\begin{longtable}{@{}P{4.54cm}P{2.59cm}P{9.79cm}@{}}
\caption{DDT\_Valve Members}\label{tab:ddt-valve}\\
\toprule \rowcolor{tblhead}\textbf{Member} & \textbf{Type} & \textbf{Description}\\ \midrule
\endfirsthead
\multicolumn{3}{@{}l}{\small\textit{Table \thetable{} continued}} \\
\toprule \rowcolor{tblhead}\textbf{Member} & \textbf{Type} & \textbf{Description}\\ \midrule
\endhead
\texttt{REQ} & BOOL & Open request from the program (automatic mode).\\
\texttt{Auto\_Perm} & BOOL & Permissive summary for automatic operation; all conditions met.\\
\texttt{Man\_Perm} & BOOL & Permissive summary for manual operation; all conditions met.\\
\texttt{HMI\_Auto} & BOOL & SCADA request to switch to Auto.\\
\texttt{HMI\_Manual} & BOOL & SCADA request to switch to Manual.\\
\texttt{HMI\_Off} & BOOL & SCADA request to switch to Off.\\
\texttt{HMI\_Lockout} & BOOL & SCADA request to switch to Lockout.\\
\texttt{CMD} & BOOL & Final open command to the output.\\
\texttt{ZSO} & BOOL & Open limit switch feedback (valve is open).\\
\texttt{ZSC} & BOOL & Closed limit switch feedback (valve is closed).\\
\texttt{HMI\_CMD} & BOOL & SCADA open/close command used in Manual mode.\\
\texttt{Interlock} & BOOL & Interlock active; the valve is forced to its safe state.\\
\texttt{Cycle\_Count} & DINT & Number of open/close cycles (maintenance counter).\\
\texttt{Reset\_Count} & BOOL & Request to reset the cycle count.\\
\texttt{Feedback\_Enabled} & BOOL & Limit switch feedback is fitted and checked (enables feedback-fault detection).\\
\texttt{Auto} & BOOL & Valve is in Auto mode.\\
\texttt{Manual} & BOOL & Valve is in Manual mode.\\
\texttt{Off} & BOOL & Valve is in Off mode.\\
\texttt{LockedOut} & BOOL & Valve is locked out.\\
\bottomrule
\end{longtable}
\section{DDT\_Motor\_VFD}
Motor driven by a variable frequency drive.

\begin{longtable}{@{}P{4.54cm}P{2.59cm}P{9.79cm}@{}}
\caption{DDT\_Motor\_VFD Members}\label{tab:ddt-motor-vfd}\\
\toprule \rowcolor{tblhead}\textbf{Member} & \textbf{Type} & \textbf{Description}\\ \midrule
\endfirsthead
\multicolumn{3}{@{}l}{\small\textit{Table \thetable{} continued}} \\
\toprule \rowcolor{tblhead}\textbf{Member} & \textbf{Type} & \textbf{Description}\\ \midrule
\endhead
\texttt{REQ} & BOOL & Start request from the program (automatic mode).\\
\texttt{INTERLOCK} & BOOL & Run interlock active; the motor is stopped.\\
\texttt{AUTO\_PERM} & BOOL & Permissive summary for automatic operation.\\
\texttt{AUTO} & BOOL & Motor is in Auto mode.\\
\texttt{MANUAL} & BOOL & Motor is in Manual mode.\\
\texttt{OFF} & BOOL & Motor is in Off mode.\\
\texttt{LOCKOUT} & BOOL & Motor is in the locked-out (service) mode \tbd.\\
\texttt{HMI\_AUTO} & BOOL & SCADA command to switch to Auto.\\
\texttt{HMI\_MANUAL} & BOOL & SCADA command to switch to Manual.\\
\texttt{HMI\_OFF} & BOOL & SCADA command to switch to Off.\\
\texttt{HMI\_LOCKOUT} & BOOL & SCADA command to switch to Lockout.\\
\texttt{CMD} & BOOL & Final motor run command to the drive.\\
\texttt{YI} & BOOL & Motor running feedback from the drive.\\
\texttt{WARNING} & BOOL & Motor warning is active.\\
\texttt{ALARM} & BOOL & Motor alarm (fault) is active.\\
\texttt{SPD\_AUTO} & REAL & Speed command used in Auto mode (from the program).\\
\texttt{SPD\_MANUAL} & REAL & Speed command used in Manual mode (from SCADA).\\
\texttt{SPD\_CMD} & REAL & Output speed command sent to the drive.\\
\texttt{RUNTIME} & REAL & Accumulated runtime of the motor \tbd (units).\\
\texttt{HMI\_RESET\_RUNTIME} & BOOL & SCADA request to reset the runtime.\\
\bottomrule
\end{longtable}
\section{DDT\_Scale}
Scale or weight transmitter interface.

\begin{longtable}{@{}P{4.54cm}P{2.59cm}P{9.79cm}@{}}
\caption{DDT\_Scale Members}\label{tab:ddt-scale}\\
\toprule \rowcolor{tblhead}\textbf{Member} & \textbf{Type} & \textbf{Description}\\ \midrule
\endfirsthead
\multicolumn{3}{@{}l}{\small\textit{Table \thetable{} continued}} \\
\toprule \rowcolor{tblhead}\textbf{Member} & \textbf{Type} & \textbf{Description}\\ \midrule
\endhead
\texttt{GrossWeight} & REAL & Gross weight reading.\\
\texttt{TareWeight} & REAL & Stored tare weight.\\
\texttt{NetWeight} & REAL & Net weight (gross minus tare).\\
\texttt{TareCmd} & BOOL & Command to tare the scale.\\
\texttt{ZeroCmd} & BOOL & Command to zero the scale.\\
\texttt{ResetScale} & BOOL & Command to reset the scale or its communications \tbd.\\
\texttt{Stable} & BOOL & Weight reading is stable (no motion).\\
\texttt{Warning} & BOOL & Scale warning is active.\\
\texttt{Alarm} & BOOL & Scale alarm is active.\\
\texttt{Fault} & BOOL & Scale or communication fault.\\
\texttt{Msg} & STRING & Status or error message text.\\
\bottomrule
\end{longtable}
\section{DDT\_Seq}
Sequence control and status.

\begin{longtable}{@{}P{4.54cm}P{2.59cm}P{9.79cm}@{}}
\caption{DDT\_Seq Members}\label{tab:ddt-seq}\\
\toprule \rowcolor{tblhead}\textbf{Member} & \textbf{Type} & \textbf{Description}\\ \midrule
\endfirsthead
\multicolumn{3}{@{}l}{\small\textit{Table \thetable{} continued}} \\
\toprule \rowcolor{tblhead}\textbf{Member} & \textbf{Type} & \textbf{Description}\\ \midrule
\endhead
\texttt{Disabled} & BOOL & Sequence is disabled.\\
\texttt{Ready} & BOOL & Sequence is ready to start (all start conditions met).\\
\texttt{Running} & BOOL & Sequence is running.\\
\texttt{Msg} & string[255] & Current status message for the operator.\\
\texttt{Start} & BOOL & Start request.\\
\texttt{Abort} & BOOL & Abort request; the sequence goes to its safe state.\\
\texttt{ShowCont} & BOOL & \tbd (likely shows the Continue prompt on SCADA).\\
\texttt{Pop\_Msg} & string[255] & Message shown in an operator pop-up \tbd.\\
\texttt{State} & DINT & Current step (state) number.\\
\texttt{Continue} & BOOL & Operator confirmation to continue past a hold point.\\
\texttt{Alarm} & DINT & Active sequence alarm code \tbd.\\
\texttt{Complete} & BOOL & Sequence has completed.\\
\texttt{Percent\_Complete} & REAL & Progress of the sequence, 0--100\%.\\
\texttt{HMI\_Auto} & BOOL & SCADA request to switch to Auto.\\
\texttt{HMI\_Manual} & BOOL & SCADA request to switch to Manual.\\
\texttt{HMI\_Off} & BOOL & SCADA request to switch to Off.\\
\texttt{Auto} & BOOL & Sequence is in Auto mode.\\
\texttt{Manual} & BOOL & Sequence is in Manual mode.\\
\texttt{Off} & BOOL & Sequence is in Off mode.\\
\texttt{Ack} & BOOL & Operator acknowledgment \tbd.\\
\texttt{HMI\_Pause} & BOOL & SCADA request to pause the sequence.\\
\texttt{Paused} & BOOL & Sequence is paused.\\
\texttt{HMI\_Topfill} & BOOL & SCADA request for a top-fill \tbd.\\
\bottomrule
\end{longtable}
\section{DDT\_Station}
Application-specific type for a cylinder fill station; used with \texttt{DFB\_Station}.

\begin{longtable}{@{}P{4.54cm}P{2.59cm}P{9.79cm}@{}}
\caption{DDT\_Station Members}\label{tab:ddt-station}\\
\toprule \rowcolor{tblhead}\textbf{Member} & \textbf{Type} & \textbf{Description}\\ \midrule
\endfirsthead
\multicolumn{3}{@{}l}{\small\textit{Table \thetable{} continued}} \\
\toprule \rowcolor{tblhead}\textbf{Member} & \textbf{Type} & \textbf{Description}\\ \midrule
\endhead
\texttt{TargetWeight} & REAL & Target fill weight.\\
\texttt{ReadyToStart} & BOOL & Station is ready to start \tbd.\\
\texttt{ReadyToFill} & BOOL & Station is ready to fill \tbd.\\
\texttt{Active} & BOOL & Station is filling.\\
\texttt{Complete} & BOOL & Station has successfully completed the fill \tbd.\\
\texttt{Deadband} & REAL & Fill deadband around the target weight \tbd.\\
\texttt{HMI\_Select} & BOOL & SCADA request to select this station.\\
\texttt{HMI\_Reset} & BOOL & SCADA request to reset the station.\\
\texttt{Selected} & BOOL & Station is selected.\\
\texttt{Venting} & BOOL & Station is venting.\\
\texttt{Vacuuming} & BOOL & Station is being evacuated.\\
\texttt{Recovering} & BOOL & Station is recovering gas \tbd.\\
\texttt{Vacuumed} & BOOL & Evacuation is complete.\\
\texttt{Press\_Balanced} & BOOL & Pressures are balanced \tbd.\\
\texttt{Disabled} & BOOL & Station is disabled.\\
\texttt{Failed} & BOOL & Station fill has failed.\\
\texttt{PT\_Valid} & BOOL & Station pressure transmitter reading is valid.\\
\texttt{PT\_Diff} & REAL & Pressure difference \tbd.\\
\texttt{InletREQ} & BOOL & Inlet valve request.\\
\texttt{InletCMD} & BOOL & Inlet valve command.\\
\texttt{SampleREQ} & BOOL & Sample valve request.\\
\texttt{SampleCMD} & BOOL & Sample valve command.\\
\texttt{VacuumREQ} & BOOL & Vacuum valve request.\\
\texttt{VacuumCMD} & BOOL & Vacuum valve command.\\
\texttt{PT} & REAL & Station pressure.\\
\texttt{Press\_Max} & REAL & Maximum allowed station pressure \tbd.\\
\texttt{TareCMD} & BOOL & Command to tare the station scale.\\
\texttt{ZeroCMD} & BOOL & Command to zero the station scale.\\
\texttt{GrossWeight} & REAL & Gross weight.\\
\texttt{TareWeight} & REAL & Tare weight.\\
\texttt{NetWeight} & REAL & Net weight.\\
\texttt{StatusEnum} & INT & Station status code (see the \texttt{StatusEnum} values in Appendix~\ref{app:function-blocks}).\\
\texttt{Msg} & string[255] & Status message for the operator.\\
\texttt{ScaleLoaded} & BOOL & A cylinder is on the scale \tbd.\\
\texttt{ScaleTared} & BOOL & Scale has been tared.\\
\texttt{Fill\_Rate} & REAL & Fill rate \tbd (units).\\
\texttt{HMI\_SP\_Max\_Vac\_Press} & REAL & Operator setpoint: maximum pressure at which the vacuum lines may be opened (permissive setpoint, Section~\ref{sec:setpoints}).\\
\texttt{TotalWeight} & REAL & Total weight \tbd.\\
\texttt{SampleWeight} & REAL & Sample weight \tbd.\\
\texttt{SamplePress} & REAL & Sample pressure \tbd.\\
\texttt{Pressurizing} & BOOL & Station is building pressure.\\
\texttt{Sampling} & BOOL & Station is sampling.\\
\texttt{HighPressWarningReq} & BOOL & Station maximum-pressure warning request \tbd.\\
\texttt{PressurizeReqd} & BOOL & Pressurizing is required \tbd.\\
\texttt{SampleReqd} & BOOL & Sampling is required \tbd.\\
\texttt{Filling} & BOOL & Station is filling.\\
\texttt{Disconnecting} & BOOL & Station is disconnecting.\\
\texttt{HMI\_Disc} & BOOL & SCADA request to disconnect.\\
\texttt{Purged} & BOOL & Station has been purged.\\
\texttt{Light} & BOOL & Station indicator light \tbd.\\
\bottomrule
\end{longtable}
